\documentclass[10pt,a4paper]{amsart}
\usepackage[margin=19mm]{geometry}
\usepackage{amsmath,amssymb,mathtools}
\usepackage[scr=rsfs,cal=euler]{mathalfa}
\usepackage{xfrac}
\usepackage[unicode,hidelinks,bookmarks=false]{hyperref}
\newcommand{\ie}{i.e.\ }
\newcommand{\cf}{cf.\ }
\newcommand{\resp}{respectively\ }
\newcommand{\Subsection}{\subsection}
\providecommand{\nobreakdash}{\nobreak\hbox{-}}
\setlength{\emergencystretch}{2em}
\multlinegap0pt
\usepackage{bm}
\usepackage{bbm}
\usepackage{dsfont}
\usepackage{xspace}
\usepackage{cancel}
\usepackage{mathtools}
\usepackage[shortlabels]{enumitem}
\usepackage{amsthm}
\makeatletter
\@ifundefined{theorem}{\newtheorem{theorem}{Theorem}}{}
\@ifundefined{prop}{\newtheorem{prop}[theorem]{Proposition}}{}
\makeatother
\title[Symmetries of extremal horizons]{Symmetries of extremal horizons}
\author{Alex Colling}
\date{Source: arXiv:2512.10200v2 (CC BY 4.0)}
\begin{document}
\maketitle

\noindent\emph{Source and licence.} Symmetries of extremal horizons. Version arXiv:2512.10200v2 (CC BY 4.0), distributed under the Creative Commons Attribution 4.0 International licence, as recorded in the arXiv licence metadata for this exact version and shown on the abstract page https://arxiv.org/abs/2512.10200v2. Full rights notice: licenses/arxiv-2512.10200-CC-BY-4.0.txt.\par\smallskip

\noindent\textit{Editorial scope.} Counted unit: the Prop whose source label is \texttt{areprop} in the article (source0.tex lines 620--626, source label \texttt{areprop}), delivered together with the article's own complete original proof of it (source0.tex lines 627--654, one \texttt{proof} environment of 28 lines). No mathematics is written, repaired or re-derived: statement and proof are the article's own text line for line. Required context retained uncounted: wave (lines 567--569); KAext (lines 573--575); I (lines 597--599). Every definition, display and label of those lines is retained unchanged. Omitted from this package and not redistributed: 1--566, 570--572, 576--596, 600--619, 655--1044. Nothing inside a retained excerpt is omitted or altered: every retained body is delivered exactly as the article's lines are written, so no omission receipt of any kind accompanies this package. Citations: the retained excerpts carry no citation command at all, so no bibliography entry of the article is transcribed and no reference is redistributed. Reference adaptation: none was needed and none was made; every reference inside the delivered unit resolves inside it, so no unresolved reference or citation is produced. Printed slips kept literal: no printed word, symbol, hypothesis or number of the counted statement or of its proof was altered. Layout and class: the article's own class is \texttt{amsart} with packages including amssymb, amsmath, hyperref, enumerate; the wrapper is the lane's pinned \texttt{amsart} editorial wrapper at 10pt on A4, which restates the theorem-like environments the retained text needs and adds the article's own macro definitions that the retained text needs (none), with extra wrapper packages (bm, bbm, dsfont, xspace, cancel, mathtools, enumitem). The delivered numbering is the unit's own. Assets: no retained span contains a figure environment, a graphics-inclusion command, a PDF-page inclusion, a table or a TikZ picture; no image file is copied into this package, the package declares no asset dependency and no image file is redistributed with it. Licence: version arXiv:2512.10200v2 of this article is distributed under the Creative Commons Attribution 4.0 International licence, as recorded in the arXiv licence metadata for this exact version and shown on the abstract page https://arxiv.org/abs/2512.10200v2. A full rights notice is delivered at licenses/arxiv-2512.10200-CC-BY-4.0.txt. Characters: every retained excerpt is pure ASCII, so the delivered engine, XeLaTeX with its default Latin Modern text font, reports no missing character.
\begin{equation} \label{wave}
    \square_{\textbf{g}}\Phi = 0.
\end{equation}

\begin{equation} \label{KAext}
      K_i(\rho,x) = \Gamma X_i(\rho,x) + \partial_i\Gamma, \hspace{1cm}   A(\rho,x) = \Gamma F(\rho,x) - \Gamma^{-1}\vert K \vert^2_{g(\rho,x)}.
\end{equation}

\begin{equation} \label{I}
    I(v,\rho) = \int_{M(v,\rho)}(2\partial_\rho \Phi + \lambda\Phi)\text{ vol}_g,
\end{equation}

\begin{prop} \label{areprop}
    Consider a solution $\Phi$ to the wave equation \eqref{wave} in a spacetime containing a doubly degenerate horizon $\mathcal{H}$.
    \begin{enumerate}[(i)]
        \item \textup{(Non-decay)} Both $I(v,0) = I_0$ and $\partial_\rho I(v,0) = I_1$ are conserved along $\mathcal{H}$, with $I$ as in \eqref{I}.
        \item \textup{(Blow-up)} If $B$ is constant and $BI_0 \neq 0$, then either $\Phi$ does not decay along $\mathcal{H}$ or $\partial_\rho^2 I(v,0)$ blows up linearly in $v$ as $v \to \infty$.
    \end{enumerate}
\end{prop}

\begin{proof}
To compute the wave operator (\ref{wave}) in the coordinates $(v,\rho,x^i)$, we require the inverse metric
\begin{equation} \label{invmet}
    \textbf{g}^{-1} = 2C\partial_v\partial_\rho - \Gamma\rho^2C^{2}A\partial_\rho\partial_\rho -2\rho CK^i\partial_\rho\partial_{x^i} + g^{ij}\partial_{x^i}\partial_{x^j}.
\end{equation}
Here $g^{ij}$ denotes the inverse of the metric $g_{ij}$ on $M(v,\rho)$, the functions $K^i$ and $A$ are given by (\ref{KAext}) and $C^{-1} = \Gamma + \rho\partial_\rho\Gamma$.
All components of (\ref{invmet}) depend on $\rho$ and $x^i$. The wave operator reads
\begin{align}\label{wavecoord}
    0 &= C^{-1}\square_{\textbf{g}}\Phi = (\text{det } g)^{-\frac 12}\partial_\mu((-\text{det }\textbf{g})^{\frac 12}\textbf{g}^{\mu\nu}\partial_\nu\Phi) = \nonumber\\&= \partial_v(2\partial_\rho\Phi + \lambda \Phi) - (\partial_\rho + \lambda)(\Gamma C\rho^2 A\partial_\rho\Phi + \rho K^i\partial_i\Phi) + \nabla_i(C^{-1}\partial^i\Phi - \rho K^i\partial_\rho \Phi).
\end{align}
We next integrate (\ref{wavecoord}) over $M(v,\rho)$. The final divergence term drops out, and in the middle term we integrate by parts to remove the derivatives $\partial_i$ acting on $\Phi$. Using $\mathcal{L}_{\partial_\rho}\text{vol}_g = \lambda\text{vol}_g$, we arrive at
\begin{equation} \label{intid}
    \partial_v I = \partial_\rho \int_{M(v,\rho)}\left(\Gamma C\rho^2A\partial_\rho \Phi -\rho (\nabla_i K^i)\Phi \right)\text{vol}_g.
\end{equation}
Since $K$ is divergence-free on the horizon, evaluating (\ref{intid}) on $\mathcal{H}$ yields
\begin{equation} \label{Ivr0}
    \partial_v I \:\overset{\mathcal{H}}{=}\:0.
\end{equation}
Hence $I = I_0$ is independent of $v$ on $\mathcal{H}$. To go further, we take a derivative of (\ref{intid}) and set $\rho = 0$,
\begin{equation}
    \partial_v \partial_\rho I \:\overset{\mathcal{H}}{=}\: 2\int_{M(v,0)}\left(A\partial_\rho\Phi - \partial_\rho(\nabla_i K^i)\Phi\right)\text{ vol}_g = 0. \label{Ivr1}
\end{equation}
The final equality holds since $A$ and $\partial_\rho(\nabla_i K^i)$ both vanish on $\mathcal{H}$. Observe that if $AI_0 \neq 0$ and $\Phi$ decays along $\mathcal{H}$ one instead finds that $\partial_v\partial_\rho I(v,0)$ approaches the constant $AI_0$, so that $\partial_\rho I(v,0)$ grows linearly in $v$. In the doubly degenerate case we must take a further derivative of (\ref{intid}),
    \begin{equation}
        \partial_v\partial_\rho^2I \:\overset{\mathcal{H}}{=}\:  \int_{M(v,0)}\left(6B \partial_\rho \Phi -3 \partial_\rho^2(\nabla_i K^i)\Phi\right)\text{vol}_g.
   \end{equation}
    If $\Phi \to 0$ as $v \to \infty$, the second term in the integrand decays as $v \to \infty$. The first term approaches $3BI_0$ provided $B$ is constant, implying the linear asymptotic growth of $\partial_\rho^2 I(v,0) \sim 3BI_0v$.
\end{proof}

\end{document}
